% ============================================================
% Public Cover Letter
%
% Copyright (c) 2026 Md Qutub Uddin Sajib.
%
% The LaTeX code and formatting may be reused under the MIT
% License. The personal, biographical, academic, employment,
% research, and other cover-letter content is not covered by
% that license and remains all rights reserved.
%
% See LICENSE for details.
%
% This source is intended for public distribution.
% Private and application-specific information is omitted.
%
% Keep this file as the canonical source for the public cover letter.
% ============================================================
%
\documentclass[11pt]{letter}

% --------------------------------------------------------------------
% Packages
% --------------------------------------------------------------------

\usepackage{microtype}
\usepackage[a4paper,
  top=1in,
  bottom=1in,
  left=1.4in,
  right=1.4in
]{geometry}
\usepackage{hyperref}
\usepackage{metalogo}
\usepackage{fontspec}
\defaultfontfeatures{Ligatures=TeX}
\setsansfont[Numbers=Lining]{Libertinus Sans}
\setmonofont[Scale=MatchLowercase]{Libertinus Mono}
\renewcommand{\familydefault}{\sfdefault}

% --------------------------------------------------------------------
% Sender information
% --------------------------------------------------------------------

\signature{%
Md Qutub Uddin Sajib\\
\href{mailto:qsajib71@gmail.com}{qsajib71@gmail.com}}

\address{}

% --------------------------------------------------------------------
% Document
% --------------------------------------------------------------------

\begin{document}

\begin{letter}{%
Hiring Committee\\
Department\\
University}

\opening{Dear Members of the Search Committee,}

% --------------------------------------------------------------------
% Opening
%
% Introduce yourself.
% State the purpose of the letter.
% Summarize your background in one or two sentences.
% --------------------------------------------------------------------

As an environmental scientist, my work combines remote sensing,
environmental modeling, geographic information systems (GIS), and
scientific computing to study environmental processes and support
evidence-based decision making.
My background includes six years of experience
in scholarly publishing, technical editing, and software development
at Mathematical Sciences Publishers. I am now seeking a postdoctoral
opportunity to apply my interdisciplinary background to collaborative
environmental research.


% --------------------------------------------------------------------
% Academic Background and Research
%
% Ph.D. in Environmental Sciences and Engineering.
% Research areas.
% --------------------------------------------------------------------
% Dissertation.
% Remote sensing.
% GIS.
% Environmental modelling.
% Scientific computing.
% Publication(s).
% --------------------------------------------------------------------

During my doctoral research, I investigated environmental processes
using satellite remote sensing, geographic information systems (GIS),
and quantitative data analysis. My work focused on the estimation of
land surface temperature from Landsat 8 imagery, including a
comparative evaluation of multiple retrieval algorithms. Through this
research, I developed a strong foundation in environmental modeling,
spatial analysis, and scientific computing using R. The work
culminated in a peer-reviewed publication and strengthened my interest
in reproducible, data-driven environmental research.


% --------------------------------------------------------------------
% Professional Experience
%
% Mathematical Sciences Publishers.
% Scholarly publishing.
% LaTeX.
% Software development.
% Workflow improvement.
% --------------------------------------------------------------------

At Mathematical Sciences Publishers, I contributed to scholarly publishing
through technical editing, production, and software development.
Collaborating with
authors, editors, mathematicians, and production staff, I contributed
to publication workflows, \LaTeX{} journal class development, and
production software improvements. This experience strengthened my
attention to detail, technical communication skills, and appreciation
for reproducible, high-quality scientific publishing.


% --------------------------------------------------------------------
% Current Research Interests
%
% Environmental data science.
% Reproducible research.
% Open-source software.
% Future research direction.
% --------------------------------------------------------------------

Today, my research interests lie at the intersection of environmental
science and scientific computing, with a particular focus on
computational approaches and open-source tools for environmental data
analysis. I am particularly interested in collaborative research
addressing complex environmental challenges through quantitative,
reproducible methods.


% --------------------------------------------------------------------
% Closing
%
% Interest in the position.
% Thank the committee.
% Availability for interview.
% --------------------------------------------------------------------

I would welcome the opportunity to contribute my interdisciplinary
background in environmental science, scientific computing, and
scholarly publishing to your research program. I look forward to
discussing how my experience and research interests could contribute
to your work.


\closing{Sincerely,}

\end{letter}

\end{document}
